\section*{Data Records}

The released corpus is provided as a JSON Lines file, \texttt{sr4all\_full.jsonl}, in the Zenodo repository (\href{https://doi.org/10.5281/zenodo.18431942}{\texttt{10.5281/zenodo.18431942}}) described in the Data Availability statement. Each line contains one systematic review encoded as a standalone JSON object, so the file can be processed sequentially without loading the complete corpus into memory. The current release contains 301,871 such review records. Full-text PDFs are not redistributed as part of the archive.

Each record is anchored by the OpenAlex work identifier stored in \texttt{id}. Where available, the record also includes the work DOI in \texttt{doi}. These identifiers link the released objects back to the underlying scholarly work and allow cross-referencing with OpenAlex metadata and citation links.

The file combines three data layers within the same review record. First, a bibliographic metadata layer is available for all records and includes fields such as \texttt{title}, \texttt{year}, \texttt{type}, \texttt{source}, \texttt{language}, \texttt{field}, \texttt{subfield}, \texttt{topics}, \texttt{keywords}, \texttt{authors}, \texttt{cited\_by\_count}, and \texttt{pdf\_url}. Second, a citation layer provides the OpenAlex reference list through \texttt{referenced\_works} together with the corresponding count in \texttt{referenced\_works\_count}. The cited works are represented by OpenAlex identifiers, yielding a consistent citation backbone for all corpus entries.

Third, records with accessible full texts and successful extraction contain an additional methodological layer derived from the extraction pipeline. Review aims are stored in \texttt{objective} and \texttt{research\_questions}. Search-related information is stored in \texttt{keywords\_used}, \texttt{exact\_boolean\_queries}, \texttt{databases\_used}, and \texttt{snowballing}. Eligibility and study-count information is stored in \texttt{inclusion\_criteria}, \texttt{exclusion\_criteria}, \texttt{n\_studies\_initial}, and \texttt{n\_studies\_final}. Temporal restrictions are represented through \texttt{year\_range}, \texttt{year\_range\_normalized}, and \texttt{year\_range\_normalization\_rule}. Missing methodological values remain \texttt{null} or are omitted when no verified value could be retained. The field \texttt{exact\_boolean\_queries} stores normalized Boolean query approximations, where available, directly in the corresponding review record rather than in a separate record type.
